% !PRINTED 202
\weberpage{202}
This remark gives us an essential limitation of the numbers $\nu$. For according to it:

\weq{13}{2n - 2 = \mu(\nu - 1) + \mu'(\nu' - 1) + \mu''(\nu'' - 1) + \cdots}

or, if we denote by $h$ the number of systems of conjugate poles, and set $n = \mu \nu = \mu' \nu' = \cdots$:

\weq{14}{2n - 2 = nh - \mu - \mu' - \mu'' - \cdots}

The numbers $\nu, \nu', \nu'' \ldots$ are at least equal to 2, hence

\[ 1 \leqq \mu \leqq \frac{n}{2}, \quad 1 \leqq \mu' \leqq \frac{n}{2}, \cdots \]

and therefore by (14)

\[ \frac{nh}{2} \leqq 2n - 2 \leqq (n-1)h, \]

or

\[ 2 \leqq h \leqq 4 - \frac{4}{n}. \]

Thus $h$ can have only one of the two values 2 or 3, and we find five ways of satisfying equation (13).

If first $h = 2$, then it follows from (14) that

\[ \mu + \mu' = 2, \quad \mu = \mu' = 1, \quad \nu = \nu' = n. \]

We therefore have:

I. \gesperrt{Cyclotomic group} or \gesperrt{cyclic group:}

\[ \nu = \nu' = n, \quad n \text{ arbitrary.} \]
\[ \mu = \mu' = 1, \]

If further $h = 3$, then it follows from (14):

\weq{15}{\mu + \mu' + \mu'' = n + 2.}

From this it is to be concluded that at least one of the numbers $\nu, \nu', \nu''$ must be equal to 2. For if they were all $\geqq 3$, we should have

\[ \mu \leqq \frac{n}{3}, \quad \mu' \leqq \frac{n}{3}, \quad \mu'' \leqq \frac{n}{3}, \]

hence $\mu + \mu' + \mu'' \leqq n$, which stands in contradiction with (15).

If then $\nu = 2$, $\mu = \dfrac{n}{2}$, it follows from (15):

\weq{16}{\mu' + \mu'' = \frac{n}{2} + 2.}

We assume first that $\nu' = 2$ as well. If we then set $\nu'' = m$, it follows from (16):

\[ \mu'' = 2, \quad n = 2m, \]
%CONT>>
